\documentclass[10pt,a4paper]{amsart}
\usepackage{amsmath,amssymb,amsthm,graphicx}
% Retained minimal preamble for the Stacks Project chapter flat.tex (More on Flatness).
% Verbatim excerpt of the chapter preamble: Xy-pic setup (native preamble lines 14-18) and the native
% theorem declarations with the mathematical macros (native preamble lines 154-217). Copied mechanically;
% nothing added, removed or redefined. The wrapper class and the amsmath/amssymb/amsthm/graphicx base
% packages are supplied by the assembler, so the class/externaldocument part of the native preamble is omitted.
% For commutative diagrams we use Xy-pic
\usepackage[all]{xy}

% We use 2cell for 2-commutative diagrams.
\xyoption{2cell}
% Theorem environments.
%
\theoremstyle{plain}
\newtheorem{theorem}[subsection]{Theorem}
\newtheorem{proposition}[subsection]{Proposition}
\newtheorem{lemma}[subsection]{Lemma}

\theoremstyle{definition}
\newtheorem{definition}[subsection]{Definition}
\newtheorem{example}[subsection]{Example}
\newtheorem{exercise}[subsection]{Exercise}
\newtheorem{situation}[subsection]{Situation}

\theoremstyle{remark}
\newtheorem{remark}[subsection]{Remark}
\newtheorem{remarks}[subsection]{Remarks}

\numberwithin{equation}{subsection}

% Macros
%
\def\lim{\mathop{\mathrm{lim}}\nolimits}
\def\colim{\mathop{\mathrm{colim}}\nolimits}
\def\Spec{\mathop{\mathrm{Spec}}}
\def\Hom{\mathop{\mathrm{Hom}}\nolimits}
\def\Ext{\mathop{\mathrm{Ext}}\nolimits}
\def\SheafHom{\mathop{\mathcal{H}\!\mathit{om}}\nolimits}
\def\SheafExt{\mathop{\mathcal{E}\!\mathit{xt}}\nolimits}
\def\Sch{\mathit{Sch}}
\def\Mor{\mathop{\mathrm{Mor}}\nolimits}
\def\Ob{\mathop{\mathrm{Ob}}\nolimits}
\def\Sh{\mathop{\mathit{Sh}}\nolimits}
\def\NL{\mathop{N\!L}\nolimits}
\def\CH{\mathop{\mathrm{CH}}\nolimits}
\def\proetale{{pro\text{-}\acute{e}tale}}
\def\etale{{\acute{e}tale}}
\def\QCoh{\mathit{QCoh}}
\def\Ker{\mathop{\mathrm{Ker}}}
\def\Im{\mathop{\mathrm{Im}}}
\def\Coker{\mathop{\mathrm{Coker}}}
\def\Coim{\mathop{\mathrm{Coim}}}

% Boxtimes
%
\DeclareMathSymbol{\boxtimes}{\mathbin}{AMSa}{"02}

%
% Macros for moduli stacks/spaces
%
\def\QCohstack{\mathcal{QC}\!\mathit{oh}}
\def\Cohstack{\mathcal{C}\!\mathit{oh}}
\def\Spacesstack{\mathcal{S}\!\mathit{paces}}
\def\Quotfunctor{\mathrm{Quot}}
\def\Hilbfunctor{\mathrm{Hilb}}
\def\Curvesstack{\mathcal{C}\!\mathit{urves}}
\def\Polarizedstack{\mathcal{P}\!\mathit{olarized}}
\def\Complexesstack{\mathcal{C}\!\mathit{omplexes}}
% \Pic is the operator that assigns to X its picard group, usage \Pic(X)
% \Picardstack_{X/B} denotes the Picard stack of X over B
% \Picardfunctor_{X/B} denotes the Picard functor of X over B
\def\Pic{\mathop{\mathrm{Pic}}\nolimits}
\def\Picardstack{\mathcal{P}\!\mathit{ic}}
\def\Picardfunctor{\mathrm{Pic}}
\def\Deformationcategory{\mathcal{D}\!\mathit{ef}}

% --- appended 2026-10-12 (Chief): the project own definitions of the reference/slogan/history
% display environments, copied mechanically from the shared Stacks preamble lines 8-12
% (\usepackage{verbatim} plus the three \newenvironment lines, which the chapter preamble keeps
% verbatim from the project preamble). Nothing added, removed or redefined: these environments
% carry no mathematical content and the project itself typesets them as comments. Without these
% three lines a counted statement whose own text uses them fails with Environment slogan undefined.
\usepackage{verbatim}
\newenvironment{reference}{\comment}{\endcomment}
\newenvironment{slogan}{\comment}{\endcomment}
\newenvironment{history}{\comment}{\endcomment}
\usepackage{cleveref}
\crefname{originaldefinition}{original definition}{original definitions}
\Crefname{originaldefinition}{Original definition}{Original definitions}
\makeatletter
\newlabel{definition-big-small-h}{{0EU1}{source}{}{sourcepointer-1}{}}
\newlabel{definition-big-small-h@cref}{{[originaldefinition][0][]0EU1}{[1][0][]source}{}{}{}}
\makeatother
\crefname{originalsection}{original section}{original sections}
\Crefname{originalsection}{Original section}{Original sections}
\makeatletter
\newlabel{example-one-generator}{{0EVG}{source}{}{sourcepointer-2}{}}
\newlabel{example-one-generator@cref}{{[originalsection][0][]0EVG}{[1][0][]source}{}{}{}}
\makeatother
\crefname{originalsection}{original section}{original sections}
\Crefname{originalsection}{Original section}{Original sections}
\makeatletter
\newlabel{example-two-generators}{{0EVH}{source}{}{sourcepointer-3}{}}
\newlabel{example-two-generators@cref}{{[originalsection][0][]0EVH}{[1][0][]source}{}{}{}}
\makeatother
\crefname{originalproposition}{original proposition}{original propositions}
\Crefname{originalproposition}{Original proposition}{Original propositions}
\makeatletter
\newlabel{proposition-check-h}{{0EVF}{source}{}{sourcepointer-4}{}}
\newlabel{proposition-check-h@cref}{{[originalproposition][0][]0EVF}{[1][0][]source}{}{}{}}
\makeatother
\crefname{originalequation}{original equation}{original equations}
\Crefname{originalequation}{Original equation}{Original equations}
\makeatletter
\newlabel{equation-almost-blow-up-square}{{0EV6}{source}{}{sourcepointer-5}{}}
\newlabel{equation-almost-blow-up-square@cref}{{[originalequation][0][]0EV6}{[1][0][]source}{}{}{}}
\makeatother
\crefname{originallemma}{original lemma}{original lemmas}
\Crefname{originallemma}{Original lemma}{Original lemmas}
\makeatletter
\newlabel{lemma-almost-blow-up-unique}{{0EVB}{source}{}{sourcepointer-6}{}}
\newlabel{lemma-almost-blow-up-unique@cref}{{[originallemma][0][]0EVB}{[1][0][]source}{}{}{}}
\makeatother
\crefname{originallemma}{original lemma}{original lemmas}
\Crefname{originallemma}{Original lemma}{Original lemmas}
\makeatletter
\newlabel{lemma-almost-blow-up-square}{{0EVA}{source}{}{sourcepointer-7}{}}
\newlabel{lemma-almost-blow-up-square@cref}{{[originallemma][0][]0EVA}{[1][0][]source}{}{}{}}
\makeatother
\crefname{originallemma}{original lemma}{original lemmas}
\Crefname{originallemma}{Original lemma}{Original lemmas}
\makeatletter
\newlabel{lemma-base-change-almost-blow-up}{{0EV7}{source}{}{sourcepointer-8}{}}
\newlabel{lemma-base-change-almost-blow-up@cref}{{[originallemma][0][]0EV7}{[1][0][]source}{}{}{}}
\makeatother
\title{A presheaf on the big h site is a sheaf iff it is a Zariski sheaf, satisfies the equalizer condition for surjective flat proper finitely presented maps to affine targets, and turns both almost blowup squares into cartesian diagrams}
\author{The Stacks Project Authors}
\date{stacks-project:flat.tex:local-snapshot-20261011; GNU FDL 1.2 or later}
\begin{document}
\maketitle
\noindent\textit{Editorial note.} All delivered mathematical text is the author's own, copied from the pinned Stacks Project chapter file flat.tex; no statement or proof step was rewritten, reordered or removed.\par\medskip
\noindent\textit{Editorial note.} The preamble is the chapter's own retained minimal preamble (sources/stacks-flat/retained-preamble.tex), a verbatim excerpt of the native Xy-pic setup and of the native theorem declarations and mathematical macros, followed by the project own definitions of the reference, slogan and history display environments (verbatim package plus three newenvironment lines) that the retained file carries at its end and that the project itself typesets as comments. No macro was invented or redefined, and the wrapper class and the amsmath/amssymb/amsthm/graphicx base packages are supplied by the assembler.\par\medskip
\noindent\textit{Editorial note.} The counted statement and proof contain no cross reference, no citation and no included file, so the delivered body is self contained.\par\medskip
\noindent\textit{Editorial note.} Canonical location of the counted proposition: Stacks tag 0EVI (sources/stacks-flat/pinned-original-source.tex line 11501).\par\medskip
\noindent\textit{Allowed standard prerequisites (editorial).}\begin{itemize}
\item Characterisation of h sheaves by the equalizer condition and the two almost blowup squares. Same chapter, Proposition labelled proposition-check-h, cited at proof lines 11521-11523. Reduces the sheaf condition to the cartesian property of an almost blow up square over an affine X.
\item The two standard almost blowup squares and their base change behaviour. Same chapter, Examples labelled example-one-generator and example-two-generators, invoked in the statement and at proof lines 11562-11564 and 11586-11588; Lemma labelled lemma-base-change-almost-blow-up, cited at proof lines 11600-11602. Provides the explicit squares to which assumptions (3) and (4) apply, and their base changes (I) and (II).
\item Uniqueness of the almost blow up square up to admissible blowups. Same chapter, Lemma labelled lemma-almost-blow-up-unique, cited at proof lines 11566-11570, and Lemma labelled lemma-almost-blow-up-square, cited at proof lines 11578-11580. Shows that whether F(X) is the fibre product does not depend on the chosen almost blow up square.
\end{itemize}\medskip
\ifdefined\hypertarget\hypertarget{sourcepointer-1}{}\fi
\noindent\textit{Source reference.} Original definition 0EU1: sources/stacks-flat/pinned-original-source.tex, source lines 10504--10504. Same-chapter pointer into the author environment definition at source line 10504; invoked at proof line None of the counted proof. It is retained only as a navigation pointer; every load-bearing step is inside the delivered span.\par\medskip
\ifdefined\hypertarget\hypertarget{sourcepointer-2}{}\fi
\noindent\textit{Source reference.} Original section 0EVG: sources/stacks-flat/pinned-original-source.tex, source lines 11474--11474. Same-chapter pointer into the author environment example at source line 11474; invoked at proof line None of the counted proof. It is retained only as a navigation pointer; every load-bearing step is inside the delivered span.\par\medskip
\ifdefined\hypertarget\hypertarget{sourcepointer-3}{}\fi
\noindent\textit{Source reference.} Original section 0EVH: sources/stacks-flat/pinned-original-source.tex, source lines 11488--11488. Same-chapter pointer into the author environment example at source line 11488; invoked at proof line 11611 of the counted proof. It is retained only as a navigation pointer; every load-bearing step is inside the delivered span.\par\medskip
\ifdefined\hypertarget\hypertarget{sourcepointer-4}{}\fi
\noindent\textit{Source reference.} Original proposition 0EVF: sources/stacks-flat/pinned-original-source.tex, source lines 11292--11292. Same-chapter pointer into the author environment proposition at source line 11292; invoked at proof line 11522 of the counted proof. It is retained only as a navigation pointer; every load-bearing step is inside the delivered span.\par\medskip
\ifdefined\hypertarget\hypertarget{sourcepointer-5}{}\fi
\noindent\textit{Source reference.} Original equation 0EV6: sources/stacks-flat/pinned-original-source.tex, source lines 10957--10957. Same-chapter pointer into the author environment equation at source line 10957; invoked at proof line 11523 of the counted proof. It is retained only as a navigation pointer; every load-bearing step is inside the delivered span.\par\medskip
\ifdefined\hypertarget\hypertarget{sourcepointer-6}{}\fi
\noindent\textit{Source reference.} Original lemma 0EVB: sources/stacks-flat/pinned-original-source.tex, source lines 11122--11122. Same-chapter pointer into the author environment lemma at source line 11122; invoked at proof line 11577 of the counted proof. It is retained only as a navigation pointer; every load-bearing step is inside the delivered span.\par\medskip
\ifdefined\hypertarget\hypertarget{sourcepointer-7}{}\fi
\noindent\textit{Source reference.} Original lemma 0EVA: sources/stacks-flat/pinned-original-source.tex, source lines 11087--11087. Same-chapter pointer into the author environment lemma at source line 11087; invoked at proof line 11600 of the counted proof. It is retained only as a navigation pointer; every load-bearing step is inside the delivered span.\par\medskip
\ifdefined\hypertarget\hypertarget{sourcepointer-8}{}\fi
\noindent\textit{Source reference.} Original lemma 0EV7: sources/stacks-flat/pinned-original-source.tex, source lines 10987--10987. Same-chapter pointer into the author environment lemma at source line 10987; invoked at proof line 11628 of the counted proof. It is retained only as a navigation pointer; every load-bearing step is inside the delivered span.\par\medskip
\begin{lemma}
\label{lemma-refine-check-h}
Let $\mathcal{F}$ be a presheaf on one of the sites $(\Sch/S)_h$
constructed in Definition \ref{definition-big-small-h}.
Then $\mathcal{F}$ is a sheaf if and only if the following
conditions are satisfied
\begin{enumerate}
\item $\mathcal{F}$ is a sheaf for the Zariski topology,
\item given a morphism $f : X \to Y$ of $(\Sch/S)_h$ with $Y$ affine
and $f$ surjective, flat, proper, and of finite presentation, then
$\mathcal{F}(Y)$ is the equalizer of the two maps
$\mathcal{F}(X) \to \mathcal{F}(X \times_Y X)$,
\item $\mathcal{F}$ turns an almost blow up square as in
Example \ref{example-one-generator} in the category $(\Sch/S)_h$
into a cartesian diagram of sets, and
\item $\mathcal{F}$ turns an almost blow up square as in
Example \ref{example-two-generators} in the category $(\Sch/S)_h$
into a cartesian diagram of sets.
\end{enumerate}
\end{lemma}

\begin{proof}
By Proposition \ref{proposition-check-h} it suffices to show that given
an almost blow up square (\ref{equation-almost-blow-up-square})
with $X$ affine in the category $(\Sch/S)_h$ the diagram
$$
\xymatrix{
\mathcal{F}(E) & \mathcal{F}(X') \ar[l] \\
\mathcal{F}(Z) \ar[u] & \mathcal{F}(X) \ar[u] \ar[l]
}
$$
is cartesian in the category of sets. The rough idea of the proof
is to dominate the morphism by other almost blowup squares to which
we can apply assumptions (3) and (4) locally.

\medskip\noindent
Suppose we have an almost blow up square
(\ref{equation-almost-blow-up-square}) in the category $(\Sch/S)_h$,
an open covering $X = \bigcup U_i$, and open coverings
$U_i \cap U_j = \bigcup U_{ijk}$ such that the diagrams
$$
\vcenter{
\xymatrix{
\mathcal{F}(E \cap b^{-1}(U_i)) & \mathcal{F}(b^{-1}(U_i)) \ar[l] \\
\mathcal{F}(Z \cap U_i) \ar[u] & \mathcal{F}(U_i) \ar[u] \ar[l]
}
}
\quad\text{and}\quad
\vcenter{
\xymatrix{
\mathcal{F}(E \cap b^{-1}(U_{ijk})) &
\mathcal{F}(b^{-1}(U_{ijk})) \ar[l] \\
\mathcal{F}(Z \cap U_{ijk}) \ar[u] &
\mathcal{F}(U_{ijk}) \ar[u] \ar[l]
}
}
$$
are cartesian, then the same is true for
$$
\xymatrix{
\mathcal{F}(E) & \mathcal{F}(X') \ar[l] \\
\mathcal{F}(Z) \ar[u] & \mathcal{F}(X) \ar[u] \ar[l]
}
$$
This follows as $\mathcal{F}$ is a sheaf in the Zariski topology.

\medskip\noindent
In particular, if we have a blow up square
(\ref{equation-almost-blow-up-square}) such that $b : X' \to X$
is a closed immersion and $Z$ is a locally principal closed
subscheme, then we see that
$\mathcal{F}(X) = \mathcal{F}(X') \times_{\mathcal{F}(E)} \mathcal{F}(Z)$.
Namely, affine locally on $X$ we obtain an almost blow up square as in (3).

\medskip\noindent
Let $Z \subset X$, $E_k \subset X'_k \to X$, $E \subset X' \to X$, and
$i_k : X' \to X'_k$ be as in the statement of
Lemma \ref{lemma-almost-blow-up-unique}. Then
$$
\xymatrix{
E \ar[d] \ar[r] & X' \ar[d] \\
E_k \ar[r] & X'_k
}
$$
is an almost blow up square of the kind discussed in the previous
paragraph. Thus
$$
\mathcal{F}(X'_k) = \mathcal{F}(X') \times_{\mathcal{F}(E)} \mathcal{F}(E_k)
$$
for $k = 1, 2$ by the result of the previous paragraph. It follows that
$$
\mathcal{F}(X) \longrightarrow
\mathcal{F}(X'_k) \times_{\mathcal{F}(E_k)} \mathcal{F}(Z)
$$
is bijective for $k = 1$ if and only if it is bijective for $k = 2$.
Thus given a closed immersion $Z \to X$ of finite presentation
with $X$ quasi-compact and quasi-separated, whether or not
$\mathcal{F}(X) = \mathcal{F}(X') \times_{\mathcal{F}(E)} \mathcal{F}(Z)$
is independent of the choice of the almost blow up square
(\ref{equation-almost-blow-up-square}) one chooses.
(Moreover, by Lemma \ref{lemma-almost-blow-up-square}
there does indeed exist an almost
blow up square for $Z \subset X$.)

\medskip\noindent
Finally, consider an affine object $X$ of $(\Sch/S)_h$
and a closed immersion $Z \to X$ of finite presentation.
We will prove the desired property for the pair $(X, Z)$
by induction on the number of generators $r$ for the ideal
defining $Z$ in $X$. If the number of generators is $\leq 2$,
then we can choose our almost blow up square as
in Example \ref{example-two-generators}
and we conclude by assumption (4).

\medskip\noindent
Induction step. Suppose $X = \Spec(A)$ and $Z = \Spec(A/(f_1, \ldots, f_r))$
with $r > 2$. Choose a blow up square
(\ref{equation-almost-blow-up-square})
for the pair $(X, Z)$. Set $Z_1 = \Spec(A/(f_1, f_2))$
and let
$$
\xymatrix{
E_1 \ar[d] \ar[r] & Y \ar[d] \\
Z_1 \ar[r] & X
}
$$
be the almost blow up square constructed in
Example \ref{example-two-generators}.
By Lemma \ref{lemma-base-change-almost-blow-up} the base changes
$$
(I)
\vcenter{
\xymatrix{
Y \times_X E \ar[r] \ar[d] & Y \times_X X' \ar[d] \\
Y \times_X Z \ar[r] & Y
}
}
\quad\text{and}\quad
(II)
\vcenter{
\xymatrix{
E \ar[r] \ar[d] & Z_1 \times_X X' \ar[d] \\
Z \ar[r] & Z_1
}
}
$$
are almost blow up squares. The ideal of $Z$ in $Z_1$ is generated
by $r - 2$ elements. The ideal of $Y \times_X Z$
is generated by the pullbacks of $f_1, \ldots, f_r$ to $Y$. Locally on $Y$
the ideal generated by $f_1, f_2$ can be generated by
one element, thus $Y \times_X Z$ is affine locally on $Y$
cut out by at most $r - 1$ elements.
By induction hypotheses and the discussion above
$$
\mathcal{F}(Y) =
\mathcal{F}(Y \times_X X') \times_{\mathcal{F}(Y \times_X E)}
\mathcal{F}(Y \times_X Z)
$$
and
$$
\mathcal{F}(Z_1) =
\mathcal{F}(Z_1 \times_X X') \times_{\mathcal{F}(E)}
\mathcal{F}(Z)
$$
By assumption (4) we have
$$
\mathcal{F}(X) =
\mathcal{F}(Y) \times_{\mathcal{F}(E_1)} \mathcal{F}(Z_1)
$$
Now suppose we have a pair $(s', t)$ with $s' \in \mathcal{F}(X')$
and $t \in \mathcal{F}(Z)$ with same restriction in $\mathcal{F}(E)$.
Then $(s'|{Z_1 \times_X X'}, t)$ are the image of a unique element
$t_1 \in \mathcal{F}(Z_1)$. Similarly,
$(s'|_{Y \times_X X'}, t|_{Y \times_X Z})$ are the image of a
unique element $s_Y \in \mathcal{F}(Y)$.
We claim that $s_Y$ and $t_1$ restrict to the same element
of $\mathcal{F}(E_1)$. This is true because the almost blow up
square
$$
\xymatrix{
E_1 \times_X E \ar[r] \ar[d] & E_1 \times_X X' \ar[d] \\
E_1 \times_X Z \ar[r] & E_1
}
$$
is the base change of almost blow up square (I) via $E_1 \to Y$ and
the base change of almost blow up square (II) via $E_1 \to Z_1$ and
because the pairs of sections used to construct $s_Y$ and $t_1$ match.
Thus by the third fibre product equality we see that there is
a unique $s \in \mathcal{F}(X)$ mapping to $s_Y$ in $\mathcal{F}(Y)$
and to $t_1$ in $\mathcal{F}(Z)$.
We omit the verification that $s$ maps to $s'$ in $\mathcal{F}(X')$
and to $t$ in $\mathcal{F}(Z)$; hint: use uniqueness of $s$ just
constructed and work affine locally.
\end{proof}

\end{document}
